% Generated draft: requires mathematical and source review.
\documentclass[11pt]{article}
\usepackage[margin=25mm]{geometry}
\usepackage{fontspec}
\setmainfont{Latin Modern Roman}
\usepackage{amsmath,amssymb,mathtools,unicode-math}
\setmathfont{Latin Modern Math}
\usepackage{enumitem,needspace,xurl,hyperref}
\setlength{\parindent}{0pt}
\setlength{\parskip}{5pt}
\setlength{\emergencystretch}{3em}
\newcommand{\N}{\mathbb{N}}
\newcommand{\Z}{\mathbb{Z}}
\newcommand{\Q}{\mathbb{Q}}
\newcommand{\R}{\mathbb{R}}
\newcommand{\C}{\mathbb{C}}
\newcommand{\F}{\mathbb{F}}
\newcommand{\A}{\mathbb{A}}
\newcommand{\PP}{\mathbb{P}}
\newcommand{\E}{\mathbb{E}}
\newcommand{\eps}{\varepsilon}
\newcommand{\dd}{\,\mathrm{d}}
\begin{document}
\section*{UnsolvedMath ID: 800702}
\textbf{Source problem code:} catalogue entry
\subsection{Definitions and mathematical statement}
For a field \(K\), let \(\operatorname{Br}(K)\) denote its Brauer group. For \(\alpha\in\operatorname{Br}(K)\), its period \(\operatorname{per}(\alpha)\) is its order in this group, and its index \(\operatorname{ind}(\alpha)\) is the degree of the central division algebra representing it. The zero class has period and index \(1\).

The supplied record uses \(\operatorname{Br.\!dim}(K)\) without defining its convention. Two possible conventions requiring distinction are
\[
 d_0(K)=\inf\left\{d\in\Z_{\ge0}:\ \operatorname{ind}(\alpha)\mid\operatorname{per}(\alpha)^d\text{ for every }\alpha\in\operatorname{Br}(K)\right\}
\]
and
\[
 d_{\mathrm{fin}}(K)=\inf\left\{d\in\Z_{\ge0}:\ \operatorname{ind}(\alpha)\mid\operatorname{per}(\alpha)^d\text{ for every finite field extension }L/K\text{ and every }\alpha\in\operatorname{Br}(L)\right\},
\]
where the infimum of the empty set is \(\infty\). The latter requires a single exponent valid uniformly over all finite extensions. The supplied excerpt does not resolve which convention is intended; these alternatives are recorded for review, without asserting their equivalence.

\textbf{Question, retaining the source notation.} For every field \(F_0\) and every extension field \(F/F_0\) finitely generated as a field extension, does
\[
 \operatorname{Br.\!dim}(F_0)<\infty
 \quad\Longrightarrow\quad
 \operatorname{Br.\!dim}(F)<\infty
\]
hold? Here finite generation means \(F=F_0(a_1,\ldots,a_m)\) for some finite list of elements, including the possibility of an algebraic extension. No characteristic restriction is stated for the general question. Once the dimension convention is fixed, the conclusion asks for a finite bound for each such field \(F\); no universal bound independent of \(F\) is specified.

\textbf{Supplied special-case context.} The source suggests first restricting the base fields. In characteristic \(p>0\), it takes \(F_0=\F_p\), the prime field, and states \(\operatorname{Br.\!dim}(F_0)=0\). For \(F\) finitely generated over this prime field, it states that transcendence degree \(1\) makes \(F\) a global field with \(\operatorname{Br.\!dim}(F)=1\), and that transcendence degree \(2\) gives \(\operatorname{Br.\!dim}(F)\le3\), citing \([\mathrm{Lieblich:PI}]\). These are supplied background assertions, not independently verified here. Their interpretation also depends on the unresolved dimension convention, and the bibliographic entry for \([\mathrm{Lieblich:PI}]\) was not supplied.
\subsection{Short English statement}
Must every finitely generated field extension of a field of finite Brauer dimension also have finite Brauer dimension? The source leaves the convention for Brauer dimension unspecified, particularly whether its period--index bound must hold uniformly over finite extensions. The registry labels the question open; this draft preserves that supplied label without independently assessing its current status.
\subsection{Sources}
\begin{itemize}
\item Local UnsolvedMath registry, permanent ID \(800702\), exact source problem code \texttt{catalogue entry}, title \emph{Open Problems on Central Simple Algebras --- Problem 12}. The record attributes the problem to Auel, Brussel, Garibaldi, and Vishne, with proposed year \(2010\), and labels it \texttt{open}. Its recorded creation and update timestamp is \texttt{2026-10-16T14:00:00Z}; this is registry metadata, not independent status verification.
\item Auel, Brussel, Garibaldi, and Vishne, \emph{Open Problems on Central Simple Algebras} (2010), Problem 12. Supplied source URL: \url{https://arxiv.org/abs/1006.3304}. The supplied provenance records access on \texttt{2026-10-16} and extraction by \texttt{arxiv\_latex}; that access is reported from the record and was not performed here.
\item The supplied statement cites \texttt{[Lieblich:PI]} for the transcendence-degree-two bound. No corresponding bibliographic entry or URL is included in the supplied record.
\end{itemize}
\end{document}
